% !TEX root = main_memo.tex

\chapter{Centered functions}\label{sectionCentered}

At this point in the paper, one could wonder if it is possible to use combinations of all the results shown already and analyze all functions in $\sC$
using their decomposition into simple functions guaranteed by \cref{JSLdecompositionlemma}.

As we saw in \cref{CBrankofPgluingofregularsequence1}, the pointed gluing of a sequence of functions is simple when the corresponding sequence of $\CB$-ranks is regular. Moreover simple functions can be decomposed in a sequence of rays at the distinguished point with regular $\CB$-ranks.
However, the latter sequence may not always be regular for continuous reducibility and the pointed gluing of this sequence of rays may not be equivalent to the simple function. % \ynote*{ajout de la référence à l'exemple}{(see \cref{infiniteWedgeOfOnes,SomeCounterExamples})}.

Scattered functions which are equivalent to a pointed gluing of some $\leq$-regular sequence play an important role in our analysis and enjoy a natural alternative definition.
We call them centered functions and we now study them in detail.

This section culminates with the results that the \bqo{} character of continuous reducibility implies that centered functions are ubiquitous.
Finally, we show that finite generation at levels of $\sC$ ensures a finite number of centered functions of any given rank, which is a crucial step in our proof.


\section{Definition and characterization}\label{ssec:centered_def}

\begin{definition}
Let $A$ be a topological space. A \emph{center} for a function $g:A\to B$ is a point $x\in A$ such that for every neighbourhood $U$ of $x$ we have $g\leq g\restr{U}$.
We say that a function $g:A\to B$ is \emph{centered} if it admits a center.\index[IdxNotion]{centered function}\index[IdxNotion]{function!centered}\index[IdxNotion]{center (of a function)}
\end{definition}

The functions $\Minimalfct{\alpha+1}$ and $\pgl\Maximalfct{\alpha}$ for all $\alpha<\omega_1$ are examples of centered functions ($\iw{0}$ is a center). More generally
saying that a sequence $(f_i)_{i\in \N}$ of functions is \emph{regular} if $\{j\in \N \mid f_i \leq f_j\}$ is infinite%\footnote{In particular, the support of $(f_i)_i$ is either empty or infinite.}
for all $i\in \N$, we have

\begin{fact}\label{Pgluingofregulariscentered}
If $(f_i)_{i\in\N}$ is a regular sequence in $\functionsonbaire$, then $\iw{0}$ is a center for $\pgl_i f_i$.
\end{fact}

\begin{proof}
Set $f:=\pgl_if_i$.
By \cref{Pgluingaslowerbound2}, it is enough to show that for every clopen neighbourhood $U$ of $\iw{0}$ and every $n\in \N$,
there exists a continuous reduction $(\sigma,\tau)$ from $f_n$ to $f$ such that $\im \sigma\subseteq U$ and $\iw{0}\notin\overline{\im(f\sigma)}$.

So let $U\ni \iw{0}$ be clopen and $m\in \N$. By regularity of $(f_i)_i$ we can find $n$ large enough such that $N_{(0)^n}\subseteq U$ and $f_m\leq f_n$, which clearly induces the desired reduction.
\end{proof}


The property of being centered is stable under continuous equivalence as a result of the following fact.

\begin{fact}\label{Centerinvariance}
Let $f,g$ be functions between topological spaces.
Suppose that $x$ is a center for $f$ and that $(\sigma,\tau)$ continuously reduces $f$ to $g$.
\begin{enumerate}
\item For every neighborhood $U$ of $\sigma(x)$, we have $f\leq g\restr{U}$.\label{Centerinvariance1}
\item If moreover $g\equiv f$, then $\sigma(x)$ is a center for $g$.\label{Centerinvariance2}
\item If $(A_i)_{i\in I}$ is an open covering of $\dom g$, then there exists $i\in I$ with $f\leq g\restr{A_i}$. \label{Centerinvariance3}
%\item If $(A_i)_{i\in I}$ is an open covering of $\dom(g)$ for some $I$, then there is an $i\in I$ with $f\leq g\restr{A_i}$.
\end{enumerate}
\end{fact}

\begin{proof}
For the first item, if $U$ is a neighbourhood of $\sigma(x)$, then by continuity of $\sigma$, $\sigma^{-1}(U)$ is a neighbourhood of $x$.
We have $f\restr{\sigma^{-1}(U)}\leq g\restr{U}$ as witnessed by $(\sigma\corestr{U},\tau)$ and $f\leq f\restr{\sigma^{-1}(U)}$ since $x$ is a center for $f$, so $f\leq g\restr{U}$.

If moreover $g\leq f$ then by transitivity $g\leq g\restr{U}$ holds. So in this case $\sigma(x)$ is a center for $g$. This gives the second point.

If now $(A_i)_i$ covers $\dom(g)$, then $\sigma(x)\in A_i$ for some $i\in I$, but as $A_i$ is open, using the first point we obtain the third one.
\end{proof}

We next establish an important interaction between the existence of centers and the property of being scattered. 

\begin{proposition}\label{scatteredhavecocenter}
Suppose that $f:A\to B$ is centered with $A$ metrizable and $B$ Hausdorff.
%Suppose that $A$ is metrizable, $B$ is Hausdorff, and $f:A\rao B$ is centered. % and continuous.
Then $f$ is scattered if and only if all centers have the same image by $f$.

Moreover when $f$ is scattered, then it is simple and any center of $f$ is mapped to its distinguished point.
\end{proposition} 
\index[IdxNotion]{scattered!function}\index[IdxNotion]{function!scattered}

\begin{proof}
Assume that $x$ is a center for a function $f$, we start by making two remarks. First, $x$ is $f$-isolated if and only if $f$ is constant. Indeed, if $f\restr{U}$ is constant on a neighborhood $U$ of $x$, then the whole $f$ is constant since we have $f\leq f\restr{U}$.

Second, for all ordinal $\alpha$ if $\CB_\alpha(f)\neq \emptyset$ then $x\in \CB_\alpha(f)$ and $x$ is a center for $f\restr{\CB_\alpha(f)}$. To see this, note that for every neighborhood $U$ of $x$ we have $f\leq f\restr{U}$ via some $(\sigma,\tau)$, hence $\sigma(\CB_\alpha(f))\subseteq \CB_\alpha(f))\cap U$ and $f\restr{\CB_\alpha(f)}\leq f\restr{\CB_\alpha(f)\cap U}$ by \cref{CBbasicsfromJSL}. So if $\CB_\alpha(f)$ is nonempty, then so is $\CB_\alpha(f)\cap U$ for all neighborhood $U$ of $x$. It follows that $x$ belongs to the closed set $\CB_\alpha(f)$ and it is a center for $f\restr{\CB_\alpha(f)}$. 

Now assume that $f$ is scattered with $\gamma=\CB(f)$, then by \cref{prop:nlc_implies_nonscattered} we have $\CB_\gamma(f)=\emptyset$. By our second remark, all centers of $f$ belong to $\CB_{\beta}(f)$ for all $\beta<\gamma$. In particular, we must have $\gamma=\alpha+1$ and all centers of $f$ belong to $\CB_\alpha(f)$. Now all centers of $f$ are centers of $f\restr{\CB_\alpha(f)}$ and they are $f\restr{\CB_\alpha(f)}$-isolated. So by our first remark, $f$ is constant on $\CB_\alpha(f)$, hence $f$ is simple and all centers have the same image.

Finally, suppose now that $x_0$ and $x_1$ are centers for $f$ such that $f(x_0)\neq f(x_1)$.
Then using our two previous remarks, we see by induction that $x_0,x_1\in\CB_\alpha(f)$ for all $\alpha$. Therefore $f$ has nonempty perfect kernel and so it is not scattered by \cref{scatterediffemptykernel}.
\end{proof}


Therefore, in particular, if $f\in\sC$ is scattered, then there is a unique point $y$ in its image such that for every center $x$ for $f$ we have $f(x)=y$.
We call this point the {\emph{\cocenter{}}} of $f$.\index[IdxNotion]{function!centered!existence of cocenter}\index[IdxNotion]{cocenter (of a centered function)}
The previous proposition has the following fundamental consequence which highlights a strong form of rigidity of centered functions with respect to continuous reducibility.

\begin{proposition}\label{Rigidityofthecocenter}
Let $f,g\in \sC$ be centered with \cocenter{} $y_f$ and $y_g$ respectively. 
Assume that $f\equiv g$ and that $(\sigma,\tau)$ is a continuous reduction from $f$ to $g$. Then the following holds:
\begin{enumerate}
\item $\tau(y_g)=y_f$,
\item for all $n\in\N$ we have $y_g\notin\overline{g\sigma(\dom \ray{f}{y_f,n})}$,
\item for all $m,n\in\N$ there is $M\geq m$ such that $\ray{f}{y_f,n}\leq \gl_{i=m}^{M}\ray{g}{y_g,i}$, % and
\item $(\ray{f}{y_f,n})_{n\in\N}$ is reducible by finite pieces to $(\ray{g}{y_g,n})_{n\in\N}$.
\end{enumerate}
\end{proposition}


\begin{proof}
First item. Let $x$ be a center for $f$. Since $f\equiv g$, $\sigma(x)$ is a center for $g$ by \cref{Centerinvariance} and so $g\sigma(x)=y_g$ by \cref{scatteredhavecocenter}.
Hence $\tau(y_g)= \tau g\sigma(x_f)=f(x_f)=y_f$ as desired. 

\smallskip

Second item: suppose not, then there is a sequence $(x_i)_{i\in \N}\subseteq\dom \ray{f}{y_f,n}$ with $g\sigma(x_i)\to y_g$, so $f(x_i)=\tau g\sigma(x_i) \to \tau(y_g)=y_f$ by continuity of $\tau$.
But by definition of rays $f(x_i)\notin \nbhd{y_f}{n+1}$ for all $i\in\N$, so $(f(x_i))_i$ cannot converge to $y_f$, a contradiction.

\smallskip

Third item: let $m,n\in\N$, use the continuity of $g$ to get $U\ni \sigma(x)$ open with $g(U)\subseteq N_{y_g\restr{m}}$.
Since $\sigma(x)$ is a center for $g$, $g\equiv g\restr{U}$ and we can choose a continuous reduction $(\sigma',\tau')$ witnessing $f\leq g\restr{U}$. %$g\leq g\restr{U}$ and we can assume that $(\sigma,\tau)$ actually reduces $f$ to $g\restr{U}$.
By the previous item, there exists $M>m$ such that $\nbhd{y_g}{M+1}$ is disjoint from $\overline{g\sigma'(\dom \ray{f}{y_f,n})}$. This shows that $(\sigma',\tau')$ restricts to the desired reduction.
\smallskip

Last item: By a recursive application of the previous item.
\end{proof}

Given a finite set $F$ of functions, we define $\pgl F=\pgl \iw{\gl F}$, and\index[IdxSymb]{ptgl@{$\pgl F$}}
we write $f\leq \FinGl{F}$\index[IdxSymb]{2@{$f\leq \FinGl{F}$}} to say that $f\leq g$ for some $g\in \FinGl{F}$.


\begin{corollary}\label{ResidualCorestrictionOfCentered}
Suppose that $f:A\to B$ in $\functionsonbaire$, $G\subseteq \sC$ finite and $f\equiv \pgl G$. Then $f$ is centered and for every open set $V\subseteq B$ which excludes its \cocenter{}, we have $f\corestr{V}\leq \FinGl{G}$.
% $y$ denote its cocenter and $V\subseteq B$ open with $y\notin V$.
%Let $f:A\to B$ be a centered function with \cocenter{} $y\in B$ and $V\subseteq B$ open with $y\notin V$.
%If $f\equiv \pgl G$ for some finite set $G\subseteq \sC$, then $f\corestr{V}\leq \FinGl{G}$. 
\end{corollary}


\begin{proof}
Suppose that $f\equiv g=\pgl G$ and let $(\sigma,\tau)$ witness $f\leq g$. By \cref{Pgluingofregulariscentered}, $g(\iw{0})=\iw{0}$ is the \cocenter{} of $g$, so $f$ is centered by \cref{Centerinvariance} and $y=\tau(\iw{0})$ is the \cocenter{} of $f$. Therefore, $(\ray{f}{y,n})_n$ is reducible by finite pieces to $(\ray{g}{\iw{0},n})_n=\iw{\gl G}$ by \cref{Rigidityofthecocenter}.
So for all $n\in\N$ we get $\ray{f}{y,n}\leq\FinGl G$ and if $V\subseteq B$ is open with $V\notni y$ then there exists $M\in\N$ with $V\subseteq \bigcup_{n<M}\ray{B}{y,n}$. Hence, it follows that $f\corestr{V}\leq \bigsqcup_{n<M}\ray{f}{y,n}\leq  \FinGl{G}$, as desired.
\end{proof}

The following theorem ties centered functions in $\sC$ to the pointed gluing operation, as promised. 
\index[IdxNotion]{operation!pointed gluing}
\index[IdxNotion]{pointed gluing operation!relation to centered functions}
\begin{theorem}\label{CenteredasPgluing}
Let $f:A\rao B$ be in $\sC$.
\begin{enumerate}
\item If $f$ is centered with \cocenter{} $y$, then $f\equiv\pgl_{n\in \N}\ray{f}{y,n}$, \label{CenteredasPgluing_def} 
\item The function $f$ is centered if and only if $f\equiv\pgl_{i\in\N} f_i$ for some monotone (or regular) sequence $(f_i)_{i\in\N}$ of functions .\label{CenteredasPgluing_mono}
\end{enumerate}
In particular, if $f$ is centered with \cocenter{} $y$, then $f$ is simple with distinguished point $y$ and $\CB(f)=(\sup_n\CB(\ray{f}{y,n}))+1$.
\end{theorem}
%\ynote{check if this result holds for a scattered $f\in\functionsonbaire$!}
%\rnote{Les fonctions $f_i$ du second point étaient supposées continues avant. Le renforcement de \cref{Pgluingaslowerbound2} nous permet plus fort. Yann: ok, mais comme on suppose $f\in\sC$, les $f_i$ sont forcément continues. Une raison pour l'hypothèse de continuité dans cette section, c'est qu'elle semblait nécessaire pour montrer l'unicité du co-centre, mais ce n'est plus le cas. À voir si ça change quelque chose.}

\begin{proof}
Assume that $f\in \sC$ is centered with \cocenter{} $y$.
For \cref{CenteredasPgluing_def}, by \cref{Pgluingofraysasupperbound} we know that $f\leq\pgl_n\ray{f}{y,n}$, so it remains to prove $f\geq\pgl_n\ray{f}{y,n}$.
To do so we apply \cref{Pgluingaslowerbound2} with $x$ a center for $f$ and let $U$ be a neighborhood of $x$: by centeredness there is a continuous reduction $(\sigma,\tau)$ from $f$ to $f\restr{U}$.
By \cref{Rigidityofthecocenter}, for all $n\in\N$ we have $f(x)=y\notin\overline{f\sigma(f^{-1}(\ray{B}{y,n}))}$, so $(\sigma,\tau)$ restricts to witness the required reductions $\ray{f}{y,n}\leq f$.

\smallskip

The reverse implication of \cref{CenteredasPgluing_mono} follows from \cref{Pgluingofregulariscentered} and \cref{Centerinvariance}.
Next we show that $f\equiv \pgl_i f_i$ for some monotone sequence $(f_i)_i$, establishing the forward implication.

By \cref{Rigidityofthecocenter}, for all $m,n\in\N$ there exists $m'>m$ such that $\ray{f}{y,n}\leq \gl_{i=m}^{m'}\ray{f}{y,i}$.
We can recursively apply this fact to get a pairwise disjoint family of finite sets $(I_n)_n$ such that, setting $f_n=\gl_{i\in I_n}\ray{f}{y,i}$,
the sequence $(f_n)_n$ is monotone, and for all $n\in\N$ we have $\ray{f}{y,n}\leq f_n$.
Moreover $(f_n)_n$ is clearly reducible by pieces to $(\ray{f}{y,n})_n$, so by a double application of \cref{Pgluingasupperbound}, we have $\pgl_nf_n\equiv\pgl_n\ray{f}{y,n}$. 

\smallskip

Finally, by \cref{scatteredhavecocenter} the distinguished point of $f$ is also its \cocenter{} $y$ and so $\CB(f)=(\sup_n \CB(\ray{f}{y,n})) +1$ by \cref{CBrankofPgluingofregularsequence2simple}.
\end{proof}


\section{Centered functions and structure of continuous reducibility}

Here we collect some important consequences of the structure of continuous reducibility on centered functions.

The first result shows that if continuous reducibility is \bqo{} up to a certain level of $\sC$, then every function at the next level is locally centered.

Recall that \bqo{} is stable under (the taking of) infinite sequences (see for instance \cite[Theorem 3.18]{wellbetterinbetween}):
if $(Q,\leq_Q)$ is \bqo{}, then the class $Q^\N$ of infinite sequences of elements of $Q$ is also \bqo{} when ordered as follows.
Given $s,t$ in $Q^\N$ say $s\leq_{Q^\N}t$\index[IdxSymb]{2@{$\leq_{Q^\N}$}} when there is a strictly increasing function $\iota:\N\to\N$ satisfying $s(x)\leq_Qt(\iota(x))$.

\index[IdxNotion]{better-quasi-order (\bqo{})}
\begin{theorem}\label{LocalCenterednessFromBQO}
For all $\alpha<\omega_1$, if $\sC_{<\alpha}$ is a \bqo{}\footnote{For those interested, $2$-\bqo{} seems to be enough. The curious reader is sent to \cite[Definition 4.1]{wellbetterinbetween}.},
then every function in $\sC_\alpha$ is locally centered.
\end{theorem}
\begin{proof}
By strong induction on $\alpha$. 

\textbf{$\alpha=0$:} The empty function, the only element of $\sC_{0}$, is trivially locally centered (although not centered).

\textbf{$\alpha$ limit:} Let $f\in \sC_{\alpha}$. Since $f$ has limit $\CB$-rank, it is locally in $\sC_{<\alpha}$. Therefore by induction hypothesis, it is locally centered.

\textbf{$\alpha$ successor:} Let $f:A\to B$ in $\sC_{\alpha}$. By the Decomposition \cref{JSLdecompositionlemma}, $f$ is locally simple,
so without loss of generality we can suppose that $f$ is simple and that $f\in\functionsonbaire$ by \cref{RepresentationforFunctions}. We let $\bar{y}$ be the distinguished point of $f$
and write $\ray{B}{i}$ and $\ray{f}{i}$ for $\ray{B}{\bar{y},i}$ and $\ray{f}{\bar{y},i}$ respectively. Note that we have $\CB(\ray{f}{i})< \alpha$ for all $i\in \N$.
Let $x\in A$, we show that it admits a neighbourhood $U$ such that $f\restr{U}$ is centered.

If there exists $s\segs x$ such that $\CB(f\restr{N_s})<\CB(f)$ then, by induction hypothesis, we are done.
Otherwise, for all $s\segs x$ we have $\CB(f\restr{N_s})=\CB(f)$.
Writing $\alpha=\beta+1$, \cref{CBbasics0}~\cref{CBbasicsfromJSL2} ensures that $N_s\cap \CB_{\beta}(f)\neq \emptyset$ for all $s\segs x$,
so $x$ belongs to the closed set $\CB_{\alpha}(f)$. Hence, we have $f(x)=\bar{y}$ and each function $f\restr{N_s}$ is simple.

Hence, for each $n$, the sequence of rays $(\ray{f}{i}\restr{N_{x\restr{n}}})_{i\in\N}$ takes value in $\sC_{<\alpha}$.
As $\sC_{<\alpha}$ is \wqo{} under continuous reducibility, we can choose by induction a non-decreasing sequence $(j_n)_n$ in $\N$
such that the sequence of functions $\rho_n=(\ray{f}{i}\restr{N_{x\restr{n}}})_{i\geq j_n}$ is regular for all $n$ (\cite[Proposition 2.9]{yann2017towardsbetter}).
Note that $m<n$ implies $\rho_m\geq_{(\sC_{<\alpha})^\N} \rho_n$, since $N_{x\restr{m}}\supseteq N_{x\restr{n}}$ implies $\ray{f}{i}\restr{N_{x\restr{m}}}\geq \ray{f}{i}\restr{N_{x\restr{n}}}$ for all $i\geq j_n$,
so $(\rho_n)_n$ is decreasing in $(\sC_{<\alpha})^\N$.
Since $\sC_{<\alpha}$ is \bqo{}, $(\sC_{<\alpha})^\N$ is \wqo{} and so there exists $m$ such that for all $n>m$ we have $\rho_m\equiv_{(\sC_{<\alpha})^\N} \rho_n$.
Define $U=N_{x\restr{m}}\setminus f^{-1}(\bigcup_{i<{j_m}}\ray{B}{i})$, we show that $f\restr{U}\equiv \pgl \rho_m$. 
Since $\rho_m$ is regular, $\pgl \rho_m$ is centered by \cref{CenteredasPgluing} and so will be $f\restr{U}$ by \cref{Centerinvariance}.

The fact that $f\restr{U}\leq \pgl \rho_m$ follows from \cref{Pgluingofraysasupperbound}. 
To show that $\pgl \rho_m \leq f\restr{U}$ using \cref{Pgluingaslowerbound2}, it is enough to show that for every $i\geq j_m$ and every $n>m$,
there exists $(\sigma,\tau)$ that continuously reduces $f^{(i)}\restr{N_{x\restr{m}}}$ to $f\restr{U}$ such that
$\im \sigma \subseteq N_{x\restr{n}}$ and $\bar{y}\notin \overline{\im f\restr{U} \sigma}$.
This is possible since $\rho_m\leq \rho_n$ and so for every $i\geq j_m$ there exists $i'\geq j_n\geq j_m$ with $f^{(i)}\restr{N_{x\restr{m}}}\leq f^{(i')}\restr{N_{x\restr{n}}}$, as desired.
\end{proof}



Pointed gluing behaves rather specifically in a finitely generated class, as pointed out by the following proposition that we use repeatedly in the sequel.


\begin{Proposition}\label{FinitegenerationandPgluing}
Let $F\subseteq \functionsonbaire$ be finite and $(f_i)_{i\in\N}$ be a sequence in $\functionsonbaire$.
\begin{enumerate}
\item If $f_i\leq\FinGl{F}$ for all $i\in\N$, then $\pgl_if_i\leq \pgl F$.
\item If for all $f\in F$ and all $i\in\N$ there is $j\geq i$ such that $f\leq f_j$, then $\pgl F\leq\pgl_if_i$.
\end{enumerate}
\end{Proposition}
\begin{proof}
For both points, we build a reduction by pieces and then use \cref{Pgluingasupperbound}.

For the first item: by hypothesis for all $n\in \N$ there is an integer $k_n$ such that $f_n\leq k_n F$. Set then $K_n=\sum_{i<n}k_i$, and $I_n=\left[K_{n}, K_{n+1}\right)$ for all $n\in\N$.
This is a family of pairwise disjoint finite subsets of $\N$ witnessing a reduction by pieces from $(f_i)_i$ to $\iw{\gl F}$.

For the second item, we build by induction a reduction by pieces.
Given $n\in\N$, suppose that we have built a family $(I_m)_{m<n}$ of pairwise disjoint finite subsets of $\N$ such that for all $m<n$ we have $\gl F\leq\gl_{i\in I_m}f_i$.
Set $j=\max(\bigcup_{m<n}I_m)+1$ and use the hypothesis to fix an injective function $\iota:F\rao[j,\infty)$ such that for all $g\in F$ we have $g\leq f_{\iota(g)}$.
Setting $I_n=\iota(F)$, which is finite since $F$ is, yields the desired reduction by pieces. %we have $\gl F\leq \gl_{i\in I_n}f_i$.
\end{proof}

\index[IdxNotion]{finitely generated set of functions}

The next theorem shows how finitely generated levels guarantee a neat representation of centered functions. Let us write $\sC_{[\lambda,\lambda+n]}= \bigcup_{i\leq n} \sC_{\lambda +i}$ for $\lambda$ limit or null and $n\in\N$.\index[IdxSymb]{Scat@{$\sC_{[\lambda,\lambda+n]}$}}

%The next theorem shows that -- except for $\Minimalfct{\lambda+1}$ -- centered functions are obtained as pointed gluing of a finite set of g from that finite generation at each Cantor\---Bendixson rank allows us to characterize centered functions in terms of these generators. 

\begin{theorem}\label{finitenessofcenteredfunctions}
Let $\lambda$ be zero or a limit ordinal and $n\in\N$.
Assume that $\sC_{[\lambda,\lambda+n]}$ is generated by some finite set $F$.
Then for every centered function $g\in\sC_{[\lambda,\lambda+n+1]}$ either $g\equiv \Minimalfct{\lambda+1}$ or there exists a non empty $G\subseteq F$ such that $g\equiv\pgl G$. 

In particular, there are finitely many centered functions up to equivalence in $\sC_{\lambda+n+1}$.
\end{theorem}
\begin{proof}
Let $g \in \sC_{[\lambda,\lambda+n+1]}$ be centered, so of successor $\CB$-rank by \cref{CenteredasPgluing}.
In particular, $g\not\equiv\Maximalfct{\lambda}$, so $\lambda<\CB(g)\leq \lambda +n+1$.

By \cref{CenteredasPgluing}, there is a $\leq$-monotone sequence $(g_i)_i$ such that $g\equiv\pgl_ig_i$ and
for every $i$ we have $\CB(g_i)<\CB(g)\leq  \lambda +n+1$ and $(\sup_i \CB(g_i))+1=\CB(g)>\lambda$. In particular, we have $\sup_i \CB(g_i) \geq \lambda$.
Assume first that $\sup_i \CB(g_i) =\lambda$. If $\lambda=0$, then $g\equiv \Minimalfct{1}= \pgl \emptyset$. Otherwise $\lambda$ is limit and $\Minimalfct{\lambda+1} \equiv \pgl_i g_i \equiv g$ by \cref{ConsequencesGeneralStructureThm}. 

Assume now that $\sup_i \CB(g_i) > \lambda$. By monotonicity there exists $j\in \N$ such that $\CB(g_i)\geq \lambda$ for all $i\geq j$, and using monotonicity again
we get $g\equiv\pgl_{i\geq j} g_i$ by two applications of \cref{Pgluingasupperbound}.
Now for every $i\geq j$, since $g_i \in \sC_{[\lambda,\lambda+n]}$ we can choose a way of writing $g_i\equiv\gl_{f\in F}n_{i,f} f$  and set $G_i=\{f\in F \mid n_{i,f}>0\}$.
We define $G=\bigcup_{i\geq j}G_i$ and show that $g\equiv\pgl G$. Note that $G$ cannot be empty, since this would imply $g_i= \emptyset$ for all $i\geq j$ and so $\sup_i \CB(g_i)=0$, a contradiction.  

By definition of $G$, $g_i\in \FinGl{G}$ for all $i\geq j$. Moreover, for all $f\in G$ we have $f\in G_i$ for some $i\geq j$, hence $f\leq g_i$ and so $f\leq g_{k}$ for all $k\geq i$ by monotonicity of $(g_i)_{i\geq j}$. Using \cref{FinitegenerationandPgluing}, we obtain $g\equiv \pgl G$.
\end{proof}

Here is the first common point between $\sC_2$ and $\sC_{\lambda+1}$ for all $\lambda$ limit that we announced.

\begin{corollary}\label{cor:CenteredSucessor}
Let $\lambda<\omega_1$ be either equal to 1 or infinite limit.
Then, up to continuous equivalence, there are two centered functions in $\sC_{\lambda+1}$: $\Minimalfct{\lambda+1}$ and $\pgl\Maximalfct{\lambda}$. Moreover, $\Minimalfct{\lambda+1}<\pgl\Maximalfct{\lambda}$ holds.
\end{corollary}

\begin{proof}
We can use \cref{finitenessofcenteredfunctions} since the hypothesis is valid thanks to \cref{LocallyConstantFunctions} for $\lambda=1$, and to \cref{JSLgeneralstructure} for $\lambda$ limit.

In case $\lambda=1$, any centered function in $\sC_2$ is equivalent to $\pgl G$ where $G$ is either $\set{\Minimalfct{1}}$, $\set{\Maximalfct{1}}$ or $\set{\Minimalfct{1},\Maximalfct{1}}$.
Therefore we get $\Minimalfct{2}$ and $\pgl\set{\Minimalfct{1},\Maximalfct{1}}\equiv\pgl\Maximalfct{1}$ by \cref{FinitegenerationandPgluing}, as wanted.

In case $\lambda$ is limit the result is straightforward since the only possible $G$ is $\set{\Maximalfct{\lambda}}$. 

For the last part, note that $\Minimalfct{\lambda+1}\leq \pgl\Maximalfct{\lambda}$ and suppose towards a contradiction that equivalence holds. If $\lambda=1$, then using \cref{Rigidityofthecocenter} we get  $\Maximalfct{1}=\id_\N \leq n \id_1=n \Minimalfct{1}$ for some $n\in\N$, a contradiction. If $\lambda$ is limit, let $\Minimalfct{\lambda+1}=\pgl_n\Minimalfct{\alpha_n+1}$ for $(\alpha_n)_n$ cofinal in $\lambda$. By \cref{Rigidityofthecocenter} again, it follows that $\Maximalfct{\lambda}\leq \gl_{n<M} \Minimalfct{\alpha_n+1}$ for some $M\in \N$, but then $\CB(\Maximalfct{\lambda})=\lambda\leq \sup_{n<M}(\alpha_n +1)<\lambda$, a contradiction again. 
\end{proof}


%We will see with \cref{OptimalityatSuccessorofLimit} that these two functions are not equivalent.

\section{Simple functions at successors of limit levels}

\index[IdxNotion]{simple function}\index[IdxNotion]{function!simple}

\begin{proposition}\label{Simpleiffcoincidenceofcocenters}
Let $f$ be in $\sC$, assume that $f=\bigsqcup_{i\in\N} f_i$ for some sequence $(f_i)_{i\in\N}$ of centered functions and set $I=\set{n\in\N}[\CB(f_n)=\sup_i\CB(f_i)]$.
\begin{enumerate}
\item $\CB(f)$ is successor if and only if $I\neq\emptyset$,
\item The $\CB$-degree of $f$ is the cardinality of the set of \cocenter{}s of the functions $f_i$ for $i$ in $I$.
\end{enumerate}
In particular, $f$ is simple if and only if $I\neq\emptyset$ and for all $n\in I$ the \cocenter{}s of the functions $f_n$ coincide with the distinguished point of $f$.
\end{proposition}

\begin{proof}
Observe first that by \cref{CBrankofclopenunion} we have $\CB(f)=\sup_i\CB(f_i)$.

To see the first item, suppose that $\CB(f)=\alpha+1$ for some $\alpha<\omega_1$. By \cref{CBbasics0}~\cref{CBbasicsfromJSL2} we have $\CB_\alpha(f)=\bigsqcup_n\CB_\alpha(f_n)$.
Since $\CB_\alpha(f)$ is non-empty, so must be $\CB_\alpha(f_n)$ for some $n$, which means that $\CB(f_n)=\alpha+1$ and $n\in I\neq\emptyset$.
If now $I\neq\emptyset$, then for any $n\in I$ we have $\CB(f_n)=\sup_i\CB(f_i)=\CB(f)$ and by \cref{CenteredasPgluing} $\CB(f_n)$ is successor, hence $\CB(f)$ is too.

\smallskip

For the second point, if $\CB(f)$ is limit then our convention is that the $\CB$-degree of $f$ is $0$, which is also the cardinality of the set of \cocenter{}s of functions in $I$, since $I=\emptyset$ by the first point. Suppose now that $\CB(f)=\alpha+1$ for some $\alpha<\omega_1$ and so $I\neq\emptyset$.
For all $n\in I$ the function $f_n$, being centered, is simple of distinguished point its \cocenter{} $y_n$ by \cref{CenteredasPgluing}.
Since by \cref{CBbasics0}~\cref{CBbasicsfromJSL2} we have $\CB_\alpha(f)=\bigsqcup_i\CB_\alpha(f_i)$ and $\CB_{\alpha}(f_i)=\emptyset$ if $i\notin I$, we actually have
$\CB_\alpha(f)=\bigsqcup_{n\in I}\CB_\alpha(f_n)$, and $f(\CB_\alpha(f))=f(\bigsqcup_{n\in I}\CB_\alpha(f_n))=\set{y_n}[n\in I]$, which proves our point.
\end{proof}


\index[IdxNotion]{Precise Structure!simple functions at successor of limit}
\begin{theorem}\label{simplefunctionslambda+1damuddafuckaz}
Let $\lambda$ be limit or 1. Assume that continuous reducibility is \bqo{} on $\sC_{<\lambda}$.
Any simple function $f\in\sC_{\lambda+1}$ is continuously equivalent to one of $\Minimalfct{\lambda+1},\Minimalfct{\lambda+1}\gl\Maximalfct{\lambda}$ or $\pgl\Maximalfct{\lambda}$.
\end{theorem}

\begin{proof}
Since we assume that $\leq$ is \bqo{} on $\sC_{<\lambda}$, by the General Structure \cref{JSLgeneralstructure} it is also \bqo{} on $\sC_{\leq\lambda}$. So by
\cref{LocalCenterednessFromBQO} we can write $f=\bigsqcup_{i\in I} f_i:A_i\to B$ with $(f_i)_{i\in I}$ a sequence of centered functions. As $f$ is simple, $f\leq \pgl\Maximalfct{\lambda}$ by \cref{Maxfunctions}, so  if $f_i\equiv\pgl\Maximalfct{\lambda}$ for some $i\in I$, then $f\equiv\pgl\Maximalfct{\lambda}$.

Since $\pgl\Maximalfct{\lambda}$ and $\Minimalfct{\lambda+1}$ are the only two centered functions in $\sC_{\lambda+1}$ by \cref{cor:CenteredSucessor},
we can suppose that for all $n\in I$ if $\CB(f_n)>\lambda$ then $f_n\equiv\Minimalfct{\lambda+1}$.
By \cref{Simpleiffcoincidenceofcocenters} we have $f_i\equiv \Minimalfct{\lambda+1}$ for at least one $i\in I$, and
the distinguished point $\bar{y}$ of $f$ is the \cocenter{} of $f_i$ for any $f_i\equiv\Minimalfct{\lambda+1}$.
In the rest of this proof, all rays are taken with respect to $\bar{y}$. 
By simplicity $\CB_\lambda(f)\subseteq f^{-1}(\set{\bar{y}})$, so $\CB(\ray{f}{n})\leq\lambda$ for all $n\in\N$.
Also we note that $\CB(\ray{f_i}{j})<\lambda$ for all $i,j\in \N$, {which follows from \cref{Rigidityofthecocenter} when $f_i\equiv\Minimalfct{\lambda+1}$.}

If $\CB(\ray{f}{n})<\lambda$ for all $n\in \N$, then $f \leq \pgl_n \ray{f}{n} \leq \Minimalfct{\lambda+1}$ by \cref{Pgluingofraysasupperbound,ConsequencesGeneralStructureThm} and so $f \equiv \Minimalfct{\lambda+1}$.

\smallskip

So we can suppose that $\CB(\ray{f}{n})=\lambda$ for some $n\in\N$. Fix $W=\ray{B}{n}$, we have $\bar{y}\notin W$, $\Maximalfct{\lambda}\leq f\corestr{W}$
by \cref{JSLgeneralstructure}, and $\Minimalfct{\lambda+1} \leq f\corestr{B\setminus W}$ by centeredness of $\Minimalfct{\lambda+1}$.
Since $W$ is clopen we have $f\corestr{B\setminus W}\gl f\corestr{W}\equiv f$ by \cref{UsefulcriterionforequivFinGl}, so $\Minimalfct{\lambda+1}\gl\Maximalfct{\lambda}\leq f$.

We show that $f\leq\Minimalfct{\lambda+1}\gl\Maximalfct{\lambda}$ also holds in that case.
To do so we are going to find a clopen partition $A=C_0\sqcup C_1$ satisfying
 $f\restr{C_0}\leq \Maximalfct{\lambda}$ and $f\restr{C_1}\leq \Minimalfct{\lambda+1}$, which by \cref{Gluingasupperbound}
 gives \(f\leq\Minimalfct{\lambda+1}\gl\Maximalfct{\lambda}\).


%\ynote*{}{Viewing $f$ outside of $f^{-1}(\bar{y})$ as partitioned into an infinite matrix $\ray{f_i}{j}$, $(i,j)\in I\times \N$, we split $f$ ``along the diagonal''. }%We split $f$ ``along the diagonal'';
We split $f$ ``along the diagonal''. For all $j\in \N$, let $g_j=\bigsqcup_{i\leq j} \ray{f_i}{j}$ and $h_{j}=\bigsqcup_{i> j} \ray{f_i}{j}$.
Consider the clopen sets $C^{j}_i=\dom (\ray{f_i}{j})=A_i\cap f^{-1}(\ray{B}{j})$ for all $i,j$. Let $C_{0}=\bigcup\{C^{j}_i\mid i > j\}$ and its complement $C_1=A\setminus C_0$. Since $\CB(\ray{f}{j}_i)<\lambda$ for all $i,j$ we have $\CB(h_j)\leq\lambda$,
so $f\restr{C_0}\leq\Maximalfct{\lambda}$ by \cref{Maxfunctions}.
For all $j$, we also have $g_j=\bigsqcup_{i\leq j}\ray{f_i}{j}$ and so $\CB(g_j)<\lambda$. Since the functions $g_j$ correspond to the rays of $f\restr{C_1}$ at $\bar{y}$, a combination of \cref{Pgluingofraysasupperbound,ConsequencesGeneralStructureThm} then yields
$f\restr{C_1}\leq\pgl_j g_j \leq\Minimalfct{\lambda+1}$.

Note that $C_0$ is open as a union of clopen sets. To see that $C_1$ is open too, note that $C_1=\bigcup_i (A_i\setminus C_0)$ and that $A_i\setminus C_0=\bigcup_{j\leq i} C^{j}_i$ is clopen for all $i$, which concludes the proof.
\end{proof}

Moreover, we show in \cref{OptimalityatSuccessorofLimit} that in fact $\Minimalfct{\lambda+1}<\Minimalfct{\lambda+1}\gl\Maximalfct{\lambda}< \pgl\Maximalfct{\lambda}$. Finally, note that
the Decomposition \cref{JSLdecompositionlemma} yields the following corollary.

\index[IdxNotion]{finitely generated set of functions}
\begin{corollary}\label{finitedegreedamuddafuckaz}
For $\lambda$ limit or 1, if continuous reducibility is \bqo{} on $\sC_{<\lambda}$, then
the set of functions in $\sC_{\lambda+1}$ that have finite degree is finitely generated by the set $\set{\Maximalfct{\lambda},\Minimalfct{\lambda+1},\pgl\Maximalfct{\lambda}}$.
\end{corollary}


